%%%%%%%%%%%%%%%%%%%%%%%%%%%%%%%%%%%%%%%%%%%%%%%%%%%%%%%%%%%%%%%%%%%%%%
% Lớp 10 | Đại số tổ hợp | Nhị thức Newton
% Dạng 0D8?3-1: Các câu hỏi lý thuyết tổng hợp
% Số câu: 10  (sinh tự động bởi tach_cau.py; từ nay thầy sửa trực tiếp file này)
%%%%%%%%%%%%%%%%%%%%%%%%%%%%%%%%%%%%%%%%%%%%%%%%%%%%%%%%%%%%%%%%%%%%%%

%%%%%%%%%%%%%%%% TRẮC NGHIỆM 4 PHƯƠNG ÁN %%%%%%%%%%%%%%%%

% nguồn: MaHoaNH_TracNghiem_Lop10
\begin{ex}%[0D8N3-1]%[ma:f4733a4c]
	Trong khai triển $(a+b)^4$, số hạng tổng quát  của khai triển là
	\choice 
		 {$\mathrm{C}_{4}^{k-1} a^k b^{5-k}$}
		{\True $\mathrm{C}_{4}^k a^{4-k} b^k$} 
		{$\mathrm{C}_{4}^{k+1}  a^{5-k} b^{k+1}$}
		{$\mathrm{C}_{4}^k a^{4-k}  b^{4-k}$}
	 \loigiai{
		Số hạng tổng quát của khai  triển $(a+b)^4$ là $\mathrm{C}_n^k a^{n-k} b^k=\mathrm{C}_{4}^k a^{4-k} b^k$.
	} 
\end{ex}

% nguồn: MaHoaNH_TracNghiem_Lop10
\begin{ex}%[0D8N3-1]%[ma:502c32e7]
	Trong khai triển $(x+1)^{2022}$  có chứa bao nhiêu số hạng?
	\choice 
		 {$2022$}
		{\True  $2023$}
		 {$2021$}
		{$2024$} 
	 \loigiai{
		Số số hạng trong khai triển $(x+1)^{2022}$ là  $2023$.
	} 
\end{ex}

% nguồn: MaHoaNH_TracNghiem_Lop10
\begin{ex}%[0D8N3-1]%[ma:58643e15]
	Trong khai triển  $(a+b)^4$ có bao nhiêu số hạng?
	 \choice
		 {$6$}
		{$3$} 
		 {\True $5$}
		{$4$} 
	 \loigiai{
		Trong khai triển $(a+b)^4$ có $4+1=5$  số hạng.
	} 
\end{ex}

% nguồn: MaHoaNH_TracNghiem_Lop10
\begin{ex}%[0D8N3-1]%[ma:6a50170a]
	Trong khai triển $(x+1)^n$ có chứa $18$ số  hạng. Vậy $n$ bằng
	\choice 
		 {$16$}
		 {\True $17$}
		 {$15$}
		 {$14$}
	 \loigiai{
		Trong khai triển $(x+1)^n$ (với $n \in \mathbb{N}$) có chứa $n+1$ số hạng, do đó $n+1=18 \Leftrightarrow  n=17$.
	} 
\end{ex}

% nguồn: MaHoaNH_TracNghiem_Lop10
\begin{ex}%[0D8N3-1]%[ma:6e360c8d]
	Trong khai triển $(2x - 5y)^4$ có bao  nhiêu số hạng?
	 \choice
		 {$6$}
		{$4$} 
		 {$7$}
		{\True  $5$}
	\loigiai{ 
		 \textbf{Nhắc lại:} Trong khai triển $(a + b)^{n}$ có $n+1$ số hạng.\\
		Như  vậy, trong khai triển $(2x - 5y)^4$ có $4 + 1 = 5$ số hạng.
	} 
\end{ex}

% nguồn: MaHoaNH_TracNghiem_Lop10
\begin{ex}%[0D8H3-1]%[ma:2cb6d87e]
	Trong khai triển  $(2x-3)^4$, số hạng tổng quát của khai triển là
	\choice 
		 {$\mathrm{C}_{4}^k 2^k 3^{4-k} \cdot x^{4-k}$}
		{\True  $\mathrm{C}_{4}^k 2^{4-k}(-3)^k \cdot x^{4-k}$}
		 {$\mathrm{C}_{4}^k 2^{4-k} 3^k \cdot x^{4-k}$}
		{$\mathrm{C}_{4}^k 2^k(-3)^{4-k} \cdot  x^{4-k}$}
	\loigiai{ 
		 Số hạng tổng quát của khai triển $(2x-3)^4$ là $\mathrm{C}_{4}^k(2x)^{4-k}(-3)^k=\mathrm{C}_{4}^k 2^{4-k}(-3)^k \cdot x^{4-k}$.
	 }
\end{ex}

% nguồn: MaHoaNH_TracNghiem_Lop10
\begin{ex}%[0D8H3-1]%[ma:9a3fbdac]
	Trong khai triển nhị thức $(x-2)^{n+2}$, $(n \in \mathbb{N})$. Có tất cả $5$ số hạng. Vậy $n$  bằng?
	\choice 
		{$3$} 
		 {\True $2$}
		{$6$} 
		 {$5$}
	 \loigiai{
		Để khai triển nhị thức $(x-2)^{n+2}$, $(n \in \mathbb{N})$ có tất cả $5$ số  hạng thì $(n+2) + 1 = 5 \Rightarrow n = 2$.
	 }
\end{ex}


%%%%%%%%%%%%%%%% TỰ LUẬN %%%%%%%%%%%%%%%%

% nguồn: MaHoaNH_TuLuan_Lop10
\begin{bt}%[0D8B3-1]%[ma:566bea20]
	Khai triển nhị thức $(2x+1)^4$. 
	 \loigiai{
		 \allowdisplaybreaks
		\begin{eqnarray*}
			 (2x+1)^4&=& \mathrm{C}_4^0(2x)^4+\mathrm{C}_4^1(2x)^3+\mathrm{C}_4^2(2x)^2+\mathrm{C}_4^3 2 x+\mathrm{C}_4^4\\
			 &=& 16x^4+32x^3+24x^2+8x+1.
		\end{eqnarray*}
	} 
\end{bt}

% nguồn: MaHoaNH_TuLuan_Lop10
\begin{bt}%[0D8B3-1]%[ma:c6206dd7]
	Khai triển biểu thức $(2x-1)^5$ theo  công thức khai triển nhị thức Newton.
	 \loigiai{
		Khai triển nhị thức Newton của $(2x-1)^5$, ta được 
		\begin{eqnarray*}
			 (2x-1)^5&=&\mathrm{C}_5^0(2x)^5+\mathrm{C}_5^1(2x)^4(-1)^1+\mathrm{C}_5^2(2x)^3(-1)^2\\
			&& +\,  \mathrm{C}_5^3(2x)^2(-1)^3+\mathrm{C}_5^4(2x)^1(-1)^4+\mathrm{C}_5^5(-1)^5\\
			&=&32x^5-80x^4+80x^3-40x^2+10x-1. 
		\end{eqnarray*}
	} 
\end{bt}

% nguồn: MaHoaNH_TuLuan_Lop10
\begin{bt}%[0D8B3-1]%[ma:f8d81a67]
	Viết khai triển và tìm số  hạng chứa $x^3$ của khai triển nhị thức Newton $(2x-3)^5$.
	 \loigiai{
		Khai triển nhị thức Newton của  $(2x-3)^5$, ta được
		\begin{eqnarray*}
			 (2x-3)^5&=&\mathrm{C}_5^0(2x)^5+\mathrm{C}_5^1(2x)^4(-3)^1+\mathrm{C}_5^2(2x)^3(-3)^2\\
			&&  +\, \mathrm{C}_5^3(2x)^2(-3)^3+\mathrm{C}_5^4(2x)^5(-3)^4+\mathrm{C}_5^5(-3)^5\\
			 &=&32x^5-240x^4+720x^3-1080x^2+810x-243.
		\end{eqnarray*}
		Từ khai  triển, ta được số hạng chứa $x^3$ là $720x^3$.
	} 
\end{bt}
